\documentclass{article}
\parindent0cm
\parskip0.5em

\renewcommand{\descriptionlabel}[1]{\hspace{\labelsep}- \texttt{#1}}

\newcommand{\entrylabel}[1]{\mbox{\textsf{\bf #1:}}\hfil}
\newenvironment{entry}
{\begin{list}{}%
  {\renewcommand{\makelabel}{\entrylabel}%
   \setlength{\labelwidth}{50pt}%
   \setlength{\leftmargin}{60pt}%
  }%
}%
{\end{list}}%

\begin{document}

\section*{Coding Standards}
\begin{enumerate}
\item[\textbf{Indentation style}] The indentation style can be characterized as the
     "indented brace". This means that the curly brace (scope delimiter) should
     be placed on the following line. It should look like this:
\begin{verbatim}
if (indicator == TRUE)
{
 doSmth();
}
\end{verbatim}
\item[\textbf{Names}] We use capitalization to indicate words within a name. To
     distinguish a local variable of a method from an attribute (global
     variable of the class), we append a trailing underline to the end of the
     attribute's name. Something like this:
\begin{verbatim}
class Foo{
 Integer varOfClass_;
}
\end{verbatim}
     Your names should be transparent, i.e. self-explanatory, and not
     incorporate non-standard abbreviations. This leads to longer names, but
     it also makes reading the code easier.
\item[\textbf{using namespace}] Directives "using namespace std" can be used only
     in function scope. It's better to use the namespace in name of variable or method. For example, it's clear to determine that variable {\tt grd::Vertex} is from namespace of grd, that is used in classes by Roberto. Also, all methods iand variables from standard library of {C++} should be mentioned with namespace {\tt std::}, like that: {\tt std::cout }. 
\item[\textbf{C++ standard libraries}] It is preferable to use methods from the
     standard C++ library and not the standard C library and also to use the
     corresponding headers. So header inclusion should look like:
\begin{verbatim}
#include <fstream>
#include <iostream>
#include <string>
\end{verbatim}
and not like:
\begin{verbatim}
#include <fstream.h>
#include <iostream.h>
#include <string.h>
\end{verbatim}
     Note that the .h versions of the C++-headers are deprecated. If they exist
     on a system, this is only for backward compatibility.
\item[\textbf{\#include}] The header file of the class should include only the
     superclasses (general classes of CFS++, such as \texttt{Vector,
     Environment, etc.}) header. All headers of standard libraries should be
     included in the source file. If you use some class in header file, you
     should provide a declaration of the class instead of including a header
     file of this class, for example in the following fashion:
     \begin{verbatim}
     class Elem;
     class Grid
     {
      Elem * listElems;
     };
     \end{verbatim}
     and not in this way:
     \begin{verbatim}
     #include "elem.hh"
     class Grid
     {
      Elem * listElems;
     };
     \end{verbatim} 
     If you have to include some files from another library, then it is better
     to do this using a flag/macro. Using this approach the code can then also
     be compiled without external classes. 
\item[\textbf{Get/Set}] We try to maintain the concept of data encapsulation and
     keep variables of a class protected. Therefore, all variables of a class
     should be declared as protected. Public access to variables can then be
     provided by appropriate Set/Get methods.
\end{enumerate}

\end{document}
